\uuid{5rZ5}
\titre{Limite}
\auteur{bodin ; hanani ; mzari}
\organisation{exo7}
\licence{CC BY-NC-SA 4.0}
\datecreate{2020-01-20}
\niveau{L1}
\module{Analyse}
\chapitre{Suite}
\sousChapitre{Autre}
\difficulte{2}
\theme{Suites, limite}
\exo7qcmid{lille-1.suites.012}

\contenu{
\begin{question}
\qclassification{121.00}
\qtags{motcle=limite}
Soit $\displaystyle u_n=\left(1-\frac{1}{n}\right)^n$ et $\displaystyle v_n=\left(1+\frac{(-1)^{n}}{n}\right)^n$. Quelles sont les bonnes réponses ?
\begin{answers}
    \good{$\displaystyle \lim _{n\to +\infty}u_n=\mathrm{e}^{-1}$}
    \bad{$\displaystyle \lim _{n\to +\infty}v_n=\mathrm{e}^{-1}$}
    \bad{La suite $(u_n)$ est divergente.}
    \good{La suite $(v_n)$ est divergente.}
\end{answers}
\begin{explanations}
On utilise le fait que si $\displaystyle \lim _{n\to +\infty}a_n=0$, alors les suites $(a_n)$ et $\ln (1+a_n)$ sont équivalentes. Ainsi
$$\ln (u_n)=n\ln \left(1-\frac{1}{n}\right)\sim n\times \frac{-1}{n}=-1.$$
Donc $(u_n)$ est convergente et sa limite est $\mathrm{e}^{-1}$. On vérifie, de même, que
$$\displaystyle \lim _{n\to +\infty}v_{2n}=\mathrm{e}\mbox{ et }\lim _{n\to +\infty}v_{2n+1}=\mathrm{e}^{-1}.$$
Donc, d'après le théorème des suites extraites, $(v_n)$ est divergente.
\end{explanations}
\end{question}
}
